\chapter{Gauge Transformations and Curvature}

In this chapter we derive the curvature of a connection on a principal bundle --- the geometric object that, for $G = U(1)$, \emph{is} the electromagnetic field strength $F$. We show that curvature is gauge-invariant (globally defined on the base), that it satisfies the Bianchi identity $\dd F = 0$ automatically, and that gauge transformations are not ad hoc symmetries but inevitable consequences of the bundle structure.

\section{The Curvature 2-Form on $P$}

\begin{definition}[Curvature]
The \textbf{curvature} of a connection $\omega$ on a principal $G$-bundle $P$ is the $\Lie{g}$-valued 2-form on $P$ defined by:
\begin{equation}
\boxed{\Omega = \dd\omega + \omega \wedge \omega.}
\end{equation}
Here $\dd$ is the exterior derivative on $P$, and $\omega \wedge \omega$ denotes the wedge product of $\Lie{g}$-valued forms combined with the Lie bracket: $(\omega \wedge \omega)(X, Y) = [\omega(X), \omega(Y)]$.
\end{definition}

This is the \textbf{structure equation} (or Cartan structure equation). The term $\omega \wedge \omega$ is quadratic in $\omega$ and involves the Lie bracket of $\Lie{g}$.

\subsection{The abelian case: $G = U(1)$}

For $G = U(1)$, the Lie algebra $\Lie{u}(1) \cong \R$ is abelian: the Lie bracket vanishes, $[\xi, \eta] = 0$ for all $\xi, \eta \in \Lie{u}(1)$. Therefore $\omega \wedge \omega = 0$, and the curvature simplifies to:

\begin{equation}
\Omega = \dd\omega.
\end{equation}

This is a dramatic simplification that is special to electromagnetism (among gauge theories). For non-abelian gauge groups (like $SU(2)$ or $SU(3)$ in the Standard Model), the quadratic term is non-zero and gives rise to self-interactions of the gauge field.

\section{The Local Curvature: The Field Strength}

\begin{definition}[Local field strength]
Given a local section $\sigma_\alpha: U_\alpha \to P$, the \textbf{local field strength} is the pullback of the curvature:
\begin{equation}
F_\alpha = \sigma_\alpha^*\Omega \in \Omega^2(U_\alpha, \Lie{g}).
\end{equation}
\end{definition}

\begin{theorem}
The local field strength is related to the local gauge potential by:
\begin{equation}
\boxed{F_\alpha = \dd A_\alpha + A_\alpha \wedge A_\alpha.}
\end{equation}
For $G = U(1)$:
\begin{equation}
F_\alpha = \dd A_\alpha.
\end{equation}
\end{theorem}

\begin{proof}
$F_\alpha = \sigma_\alpha^*\Omega = \sigma_\alpha^*(\dd\omega + \omega \wedge \omega) = \dd(\sigma_\alpha^*\omega) + (\sigma_\alpha^*\omega)\wedge(\sigma_\alpha^*\omega) = \dd A_\alpha + A_\alpha \wedge A_\alpha$.
\end{proof}

For electromagnetism, $F = \dd A$ is exactly the relation between the field strength 2-form and the gauge potential 1-form that we wrote down in Chapter~5. But now we understand \emph{where it comes from}: $F$ is the curvature of the $U(1)$ connection, pulled back to spacetime via a local section.

\section{Gauge Invariance of the Curvature}

\begin{theorem}
On overlapping patches $U_\alpha \cap U_\beta$, the local field strengths agree:
\begin{equation}
F_\beta = g_{\alpha\beta}^{-1}\,F_\alpha\,g_{\alpha\beta}.
\end{equation}
For $G = U(1)$ (abelian), this gives simply $F_\beta = F_\alpha$.
\end{theorem}

\begin{proof}
The curvature $\Omega$ is a tensorial 2-form on $P$ (it is $\Ad$-equivariant and horizontal). The pullback through different sections gives gauge-related but equivalent expressions. For $U(1)$, since $\Ad_g = \id$, we get $F_\beta = F_\alpha$ on overlaps.
\end{proof}

This is the geometric explanation of gauge invariance: the field strength $F$ is the local representative of a \emph{globally defined, gauge-invariant} 2-form on $M$. Unlike $A$, which is section-dependent, $F$ patches together consistently across all of $M$.

\begin{remark}
This explains the asymmetry between $A$ and $F$ that was puzzling in Chapter~5:
\begin{itemize}
\item $A$ is a local representative of a connection --- it depends on the choice of section (gauge).
\item $F$ is a local representative of the curvature --- it is independent of the section and globally defined.
\end{itemize}
The connection is not a tensor; the curvature is.
\end{remark}

\section{Gauge Transformations Revisited}

\subsection{Gauge transformations as bundle automorphisms}

\begin{definition}[Gauge transformation]
A \textbf{gauge transformation} is a bundle automorphism $\varphi: P \to P$ that:
\begin{enumerate}
\item covers the identity: $\pi \circ \varphi = \pi$ (it does not move points on the base),
\item commutes with the $G$-action: $\varphi(u \cdot g) = \varphi(u) \cdot g$.
\end{enumerate}
\end{definition}

A gauge transformation is completely determined by a smooth function $g: M \to G$ via $\varphi(u) = u \cdot g(\pi(u))$. The group of all gauge transformations is $\mathcal{G} = \Ck(M, G)$.

\subsection{Effect on the connection and curvature}

Under a gauge transformation determined by $g: M \to G$:
\begin{align}
A &\to g^{-1}Ag + g^{-1}\dd g, \label{eq:gauge_A} \\
F &\to g^{-1}Fg. \label{eq:gauge_F}
\end{align}

For $U(1)$, writing $g(p) = e^{i\chi(p)}$:
\begin{align}
A &\to A + \dd\chi, \\
F &\to F. \qquad \text{(invariant)}
\end{align}

The gauge potential transforms inhomogeneously (it is not a tensor); the field strength transforms homogeneously (it is a tensor). For abelian groups, the field strength is strictly invariant.

\section{The Bianchi Identity}

\begin{theorem}[Bianchi identity]
The curvature satisfies
\begin{equation}
\boxed{\dd_\omega \Omega = 0,}
\end{equation}
where $\dd_\omega = \dd + [\omega, \cdot\,]$ is the covariant exterior derivative.
\end{theorem}

\begin{proof}
Applying $\dd$ to $\Omega = \dd\omega + \omega\wedge\omega$:
\begin{align}
\dd\Omega &= \dd^2\omega + \dd\omega\wedge\omega - \omega\wedge\dd\omega \\
&= (\Omega - \omega\wedge\omega)\wedge\omega - \omega\wedge(\Omega - \omega\wedge\omega) \\
&= \Omega\wedge\omega - \omega\wedge\Omega - \omega\wedge\omega\wedge\omega + \omega\wedge\omega\wedge\omega \\
&= \Omega\wedge\omega - \omega\wedge\Omega.
\end{align}
Hence $\dd\Omega + \omega\wedge\Omega - \Omega\wedge\omega = 0$, which is $\dd_\omega\Omega = 0$.
\end{proof}

For $U(1)$, the covariant exterior derivative reduces to the ordinary exterior derivative (the bracket terms vanish), and the Bianchi identity becomes:
\begin{equation}
\dd F = 0.
\end{equation}

This is the homogeneous Maxwell equation, now revealed as a purely geometric identity --- the Bianchi identity for the curvature of a $U(1)$ connection. It is \emph{not} a dynamical equation; it is an automatic consequence of the fact that $F$ is a curvature.

\section{Flat Connections}

\begin{definition}[Flat connection]
A connection is \textbf{flat} if its curvature vanishes: $\Omega = 0$, i.e.\ $F = 0$.
\end{definition}

A flat connection is locally trivial (by the Frobenius theorem, the horizontal distribution is integrable), but it may be globally non-trivial if the base manifold has non-trivial topology. This is precisely the situation in the Aharonov--Bohm effect, which we study in the next chapter.

\begin{proposition}
Flat $U(1)$-connections on $M$ are classified (up to gauge equivalence) by homomorphisms $\pi_1(M) \to U(1)$, i.e.\ by $\Hom(\pi_1(M), U(1))$.
\end{proposition}

If $\pi_1(M)$ is trivial (simply connected base), every flat connection is gauge-equivalent to the trivial connection. If $\pi_1(M) \neq 0$, there exist physically inequivalent flat connections, distinguished by their holonomy.

\section{Summary}

\begin{table}[h]
\centering
\renewcommand{\arraystretch}{1.4}
\begin{tabular}{lll}
\toprule
\textbf{Object} & \textbf{Expression} & \textbf{Physical meaning} \\
\midrule
Curvature on $P$ & $\Omega = \dd\omega + \omega\wedge\omega$ & Intrinsic field strength \\
Local field strength & $F = \dd A$ (for $U(1)$) & EM field on spacetime \\
Gauge transformation of $A$ & $A \to A + \dd\chi$ & Change of local section \\
Gauge invariance of $F$ & $F \to F$ & $F$ is globally defined \\
Bianchi identity & $\dd F = 0$ & Homogeneous Maxwell (automatic) \\
Flat connection & $F = 0$ & Field-free, but holonomy may be non-trivial \\
\bottomrule
\end{tabular}
\end{table}

The curvature of the connection is the field strength. The Bianchi identity is the homogeneous Maxwell equation. Gauge invariance of $F$ is a theorem about curvature tensors. In the next chapter, we explore what happens when the curvature vanishes but the connection is still physically detectable: the Aharonov--Bohm effect.
